\section{Conclusiones}
\label{sec:conclusiones}

El proyecto alcanzó los objetivos planteados, demostrando la viabilidad de implementar un sistema de gestión y seguimiento de tesis utilizando Laravel con Inertia.js y la metodología UWE para el modelado.

\begin{itemize}
    \item Se logró implementar la totalidad de los 8 casos de uso definidos en la etapa de levantamiento de requisitos, cubriendo desde la gestión de usuarios hasta la generación de reportes y estadísticas.
    \item La metodología UWE (UML-based Web Engineering) proporcionó un marco estructurado para el modelado del sistema, permitiendo documentar de manera sistemática los modelos de contenido, navegación, presentación y procesos mediante diagramas UML estandarizados.
    \item La combinación de Scrum como marco ágil con sprints semanales permitió un avance constante y medible, completando 20 historias de usuario en 8 sprints con un esfuerzo total de 118 horas.
    \item La arquitectura monolítica con Laravel e Inertia.js demostró ser adecuada para el alcance del proyecto, eliminando la complejidad de una API REST independiente y facilitando el desarrollo rápido con un único servidor de aplicaciones.
    \item La integración con Stripe para pagos electrónicos se implementó exitosamente, soportando las modalidades de contado y crédito en 2 cuotas, con registro persistente de transacciones.
    \item El uso de Eloquent ORM y migraciones de Laravel facilitó la gestión del esquema de base de datos en PostgreSQL, con 13 tablas principales y relaciones correctamente definidas.
\end{itemize}

\endinput
